\chapter*{PHỤ LỤC}
\addcontentsline{toc}{chapter}{PHỤ LỤC}
\markboth{PHỤ LỤC}{}

Phụ lục liệt kê toàn bộ tài nguyên công khai phục vụ việc tái lập (reproducibility) kết quả đề
tài. Tất cả các URL dưới đây có thể truy cập công khai tại thời điểm nộp đề tài.

\section*{Phụ lục A: Mã nguồn (GitHub)}
\addcontentsline{toc}{section}{Phụ lục A: Mã nguồn (GitHub)}
\label{app:github}

\begin{itemize}[leftmargin=1.5cm]
\item \textbf{Repository chính:}\\
\url{https://github.com/huynhphtloi/rgbd-object-understanding}
\item \textbf{Cấu trúc thư mục quan trọng:}
\begin{itemize}
\item \texttt{config/} --- tệp cấu hình YAML (dữ liệu, huấn luyện, demo, tri thức kích thước)
\item \texttt{src/data/} --- tải và chuyển đổi OCID, COCO; sinh cảnh tổng hợp
\item \texttt{src/segmentation/} --- huấn luyện và suy luận YOLO-seg
\item \texttt{src/depth/} --- chiếu ngược, đám mây điểm, tách mặt phẳng, MiDaS
\item \texttt{src/recognition/} --- nhận diện từ vựng mở bằng CLIP kèm cơ chế từ chối
\item \texttt{src/knowledge/} --- bảng kích thước tiên nghiệm và hiệu chỉnh tỉ lệ
\item \texttt{src/reasoning/} --- luật quan hệ và dựng đồ thị cảnh
\item \texttt{src/evaluation/} --- sáu kịch bản đánh giá dùng trong Chương~\ref{ch:ket_qua}
\item \texttt{outputs/metrics/} --- tệp JSON kết quả gốc của mọi con số đã báo cáo
\item \texttt{tests/} --- 74 kiểm thử tự động
\item \texttt{notebooks/} --- notebook Colab chạy đầu-cuối
\item \texttt{report/} --- mã nguồn \LaTeX{} của đề tài này
\end{itemize}
\end{itemize}

\section*{Phụ lục B: Bộ dữ liệu và trọng số}
\addcontentsline{toc}{section}{Phụ lục B: Bộ dữ liệu và trọng số}
\label{app:dataset}

\begin{itemize}[leftmargin=1.5cm]
\item \textbf{OCID --- Object Clutter Indoor Dataset:}\\
{\small\sloppy\url{https://www.acin.tuwien.ac.at/en/vision-for-robotics/software-tools/object-clutter-indoor-dataset/}\par}
\begin{itemize}
\item Quy mô sử dụng: tập kiểm thử 26 chuỗi cảnh, 312 khung hình
\item Đơn vị độ sâu: milimet (hệ số quy đổi $0{,}001$ về mét)
\item Ngày truy cập: 04/08/2026
\end{itemize}

\item \textbf{COCO val2017:} \url{https://cocodataset.org/}
\begin{itemize}
\item Quy mô sử dụng: 200 ảnh đầu tiên theo thứ tự tất định
\item Nhãn sử dụng: mặt nạ instance và tên lớp do người gán
\item Ngày truy cập: 04/08/2026
\end{itemize}

\item \textbf{Trọng số phân vùng đã huấn luyện:}\\
{\small\sloppy\url{https://huggingface.co/loihuynh/rgbd-object-understanding-yolo-seg}\par}
\begin{itemize}
\item Tệp: \texttt{best.pt} --- kiến trúc YOLOv8n-seg, một lớp \texttt{object}
\item Huấn luyện: 100 epoch, kích thước ảnh 640, batch 8, hạt giống 42
\item Mã băm SHA-256:\\
{\footnotesize\texttt{152a8923e81d14d3a90a9204c4be17bc}\\
\texttt{76de5c1a5caf1845cb423a7c0cafd771}\par}
\end{itemize}
\end{itemize}

\hopnhanmanh[Lưu ý khi tái lập]{Nếu thiếu tệp \texttt{best.pt}, hệ thống vẫn chạy được nhưng sẽ
tự động lùi về mô hình COCO \texttt{yolov8m-seg} dùng theo kiểu class-agnostic. Đó \emph{không}
phải hệ thống đã được đánh giá và sẽ không tái lập được các con số trong báo cáo. Cần đối chiếu mã băm SHA-256 trước khi trích dẫn bất kỳ kết quả nào trong báo cáo.}

\section*{Phụ lục C: Hướng dẫn tái lập}
\addcontentsline{toc}{section}{Phụ lục C: Hướng dẫn tái lập}
\label{app:repro}

\begin{lstlisting}[language=bash,caption={Các bước tái lập toàn bộ kết quả của đề tài}]
# 1. Cai dat moi truong
python3 -m venv .venv && source .venv/bin/activate
pip install -r requirements.txt

# 2. Tai trong so da huan luyen va doi chieu ma bam
python3 -c "
from huggingface_hub import hf_hub_download
import shutil, pathlib
dst = pathlib.Path('outputs/models/rgb_yolo_seg_baseline/weights')
dst.mkdir(parents=True, exist_ok=True)
src = hf_hub_download('loihuynh/rgbd-object-understanding-yolo-seg', 'best.pt')
shutil.copy(src, dst / 'best.pt')"
shasum -a 256 outputs/models/rgb_yolo_seg_baseline/weights/best.pt

# 3. Kiem thu nhanh - khong can tai du lieu
python3 -m pytest                                    # 74 tests

# 4. Chuan bi du lieu
python3 -m src.data.download_ocid   --out data/raw/ocid
python3 -m src.data.split_by_scene  --root data/raw/ocid
python3 -m src.data.fetch_coco_gt   --out data/samples/coco_val --limit 200

# 5. Sau thi nghiem danh gia (theo dung thu tu trong Chuong 4)
python3 -m src.evaluation.eval_segmentation
python3 -m src.evaluation.eval_counting              --limit 200
python3 -m src.evaluation.eval_identification_coco   --images data/samples/coco_val
python3 -m src.evaluation.eval_scale_control         --synthetic
python3 -m src.evaluation.eval_measurement_rgbd      --root data/raw/ocid --limit 50
python3 -m src.evaluation.eval_relations_synthetic
\end{lstlisting}

\textbf{Môi trường đã dùng để tạo ra kết quả:} Python 3.11; các thư viện chính được ghim phiên
bản chính xác trong \texttt{requirements.txt} (PyTorch 2.8.0, Ultralytics 8.4.67, open\_clip\_%
torch 3.3.0, Open3D 0.18.0, NumPy 2.0.2); hạt giống ngẫu nhiên 42 cho cả bước chia dữ liệu lẫn
bước huấn luyện. Mọi tệp JSON kết quả gốc đều được lưu trong thư mục
\texttt{outputs/metrics/} của kho mã nguồn, để đối chiếu trực tiếp với các bảng số liệu ở
Chương~\ref{ch:ket_qua}.
